\section{Related Work}
\label{sec:related_work}

This section proceeds from the imaging modality and its dimensionality problem (Section~\ref{sec:related_hsi}), to the band-selection methods available for conventional cubes (Section~\ref{sec:related_band}), to the deep-learning and attribution tools on which the framework builds (Sections~\ref{sec:related_dl} and~\ref{sec:related_explain}), and closes with the gap the framework addresses (Section~\ref{sec:related_summary}).

\subsection{Multi-Excitation Fluorescence Imaging and the Dimensionality Problem}
\label{sec:related_hsi}

Hyperspectral imaging systems have evolved from early pushbroom scanners requiring platform motion~\cite{goetz1985imaging} to contemporary snapshot sensors capable of real-time acquisition~\cite{hagen2013review}.
While reflectance-based hyperspectral imaging remains dominant for remote sensing applications, fluorescence hyperspectral imaging has gained traction in laboratory and close-range settings where controlled excitation is feasible~\cite{mishra2017close}.
The extension to multiple excitation wavelengths---yielding the EEM structure that characterizes ME-HSI---was pioneered in analytical chemistry for fluorescence spectroscopy of solutions~\cite{murphy2013fluorescence} and has been adapted for imaging applications where spatially resolved EEM data provides additional discrimination power~\cite{lavine2001exploratory}.
Such systems have proven particularly valuable for biological specimens, including lichens and plant tissues, where autofluorescence signatures encode information about molecular composition that reflectance imaging cannot capture~\cite{croce2019autofluorescence}.
Recent reviews of multi- and hyperspectral autofluorescence imaging document a growing clinical and biological application base~\cite{campbell2023beyondoncology}, and multi-excitation autofluorescence hyperspectral imaging has been demonstrated for the exploration of biological samples~\cite{ghaffari2019multiexcitation}.

The 4D data structure of ME-HSI presents unique challenges.
Unlike 3D hyperspectral cubes where spectral bands sample a single physical phenomenon (reflectance or emission at fixed excitation), ME-HSI data exhibits complex cross-excitation correlations arising from energy transfer between fluorophores, concentration-dependent quenching, and overlapping emission spectra~\cite{lakowicz2006principles}.
Additionally, Rayleigh and Raman scattering at wavelengths near the excitation produce artifacts that must be removed, resulting in variable numbers of valid emission bands per excitation~\cite{murphy2013fluorescence}.
While hardware advances have made 4D acquisition practical, analytical methods for handling the resulting data dimensionality have lagged behind~\cite{mishra2017close}.

\subsection{Band Selection for Hyperspectral Cubes}
\label{sec:related_band}

Band selection methods aim to identify a subset of spectral bands that retain discriminative information while reducing dimensionality~\cite{sun2019hyperspectral,sawant2021survey}.
The main families, their specific contribution and their limitation for ME-HSI data are as follows.
\begin{enumerate}
\item \textbf{Filter methods} rank bands by statistics computed without reference to any downstream task: variance~\cite{chang2000constrained}, entropy or mutual information~\cite{guo2006band}, and inter-band correlation~\cite{martinez2006comparison}.
  Contribution: negligible cost and no labels.
  Limitation: no notion of discriminative value (a high-variance band may capture noise) and independent treatment of bands, so excitation-emission combinations are never assessed jointly.
\item \textbf{Wrapper methods} score band subsets by the performance of a classifier trained on them: sequential forward selection~\cite{pudil1994floating} and genetic or particle-swarm search~\cite{pal2010feature}.
  Contribution: task-relevant subsets.
  Limitation: labels are required and the cost grows combinatorially with the number of bands; exhaustive evaluation of a 192-band ME-HSI grid is intractable.
\item \textbf{Sparse and embedded supervised methods} fold selection into model training through $L_1$ regularization~\cite{tibshirani1996regression} or attention layers trained for classification~\cite{lorenzo2020attention}.
  Contribution: selection and prediction optimized jointly.
  Limitation: labeled training data are required, limiting applicability when ground truth is scarce or expensive to obtain.
\item \textbf{Attention--reconstruction networks} learn band weights jointly with a reconstruction network and read the selection off the weights: BS-Nets~\cite{cai2020bsnets}, attention-based autoencoders~\cite{dou2021attentionae}, DARecNet-BS~\cite{roy2021darecnet} and TAttMSRecNet~\cite{nandi2023tattmsrecnet}.
  Contribution: label-free, end-to-end band weighting.
  Limitation: a single attention vector over a single spectral axis presumes one rectangular cube, and the selection is fixed at training time.
\item \textbf{Concrete-autoencoder selectors} make subset selection differentiable through a relaxed categorical layer~\cite{balin2019concrete,sun2022concrete}.
  Contribution: exact-$K$ unsupervised selection with a reconstruction objective.
  Limitation: $K$ is an architectural constant and the selector layer again assumes a single spectral axis.
\item \textbf{Reinforcement-learning selectors} cast band selection as a sequential decision process~\cite{mou2022reinforcement,guo2025twostage}.
  Contribution: explicit search with a learned policy.
  Limitation: the reward is a classification score or a per-step proxy, and the search cost grows with the number of candidate bands.
\item \textbf{Reconstruction-based combinatorial selection} chooses bands that best reconstruct their neighbours~\cite{wang2020onr}.
  Contribution: an unsupervised, interpretable objective.
  Limitation: band-to-band linear reconstruction of a 3D cube, without cross-excitation structure.
\item \textbf{Wavelength selection for fluorescence imaging} has been addressed for autofluorescence tissue imaging~\cite{asfour2018}, on the same modality family as the present work.
  Contribution: task-driven selection of acquisition wavelengths.
  Limitation: supervised and single-excitation.
\end{enumerate}

Critically, all of these methods assume a 3D data structure in which bands can be evaluated independently or along a single spectral axis.
For ME-HSI, where information lies in excitation-emission \emph{pairs} and cross-excitation patterns may be discriminative, methods that explicitly model the 4D structure are needed.
Recent work on tensor decomposition for EEM data~\cite{murphy2013fluorescence} and multi-way analysis~\cite{bro1997parafac} addresses spectroscopic rather than imaging data, and the closest imaging work, multi-excitation autofluorescence hyperspectral imaging of biological samples~\cite{ghaffari2019multiexcitation}, analyses the data by chemometric decomposition rather than by band selection.
The present framework bridges this gap by extending unsupervised deep band selection to the 4D imaging case.

\subsection{Deep Learning and Autoencoders for Hyperspectral Data}
\label{sec:related_dl}

Deep learning has substantially advanced hyperspectral image analysis, with convolutional neural networks achieving state-of-the-art results in classification, segmentation, and anomaly detection~\cite{li2019deep,paoletti2019deep}; recent surveys trace the evolution from shallow to deep feature extraction~\cite{rasti2020featureextraction,ahmad2022survey}.
Early approaches applied 2D CNNs to spectral vectors treated as 1D signals~\cite{hu2015deep}, while subsequent work exploited spatial context through 3D convolutions that jointly process spatial and spectral dimensions~\cite{chen2016deep}.
Architectures such as spectral-spatial residual networks~\cite{zhong2017spectral} and HybridSN~\cite{roy2020hybridsn} combine 3D convolutions for spectral-spatial feature extraction with 2D convolutions for spatial refinement, achieving classification accuracies exceeding 95\% on standard benchmarks.

For unsupervised analysis, autoencoders have been applied to hyperspectral data for compression~\cite{chen2014deep}, denoising~\cite{vincent2008extracting}, and feature extraction~\cite{mou2017unsupervised}.
Variational autoencoders (VAEs) provide probabilistic latent representations that enable uncertainty quantification~\cite{kingma2013auto}, while sparse autoencoders encourage disentangled features through sparsity constraints~\cite{ng2011sparse}.
Stacked convolutional autoencoders~\cite{masci2011stacked} have demonstrated the ability to learn hierarchical features from spectral data without supervision.
These works focus on learning representations for downstream classification; in contrast, the present framework uses the learned representation as the basis for wavelength attribution.

Most deep learning approaches for hyperspectral data assume a 3D input structure.
While 3D convolutions naturally extend to 4D data by adding an excitation dimension, the variable number of emission bands per excitation after Rayleigh cutoff complicates a monolithic architecture.
The parallel-branch design introduced here addresses this limitation, enabling end-to-end learning on 4D ME-HSI data with heterogeneous spectral dimensions.

\subsection{Perturbation-Based Attribution}
\label{sec:related_explain}

Explaining neural network predictions has become crucial for scientific applications where understanding \emph{why} a model makes decisions is as important as accuracy~\cite{gilpin2018explaining}.
Post-hoc explanation methods include saliency maps that compute gradients of outputs with respect to inputs~\cite{simonyan2014deep}, Grad-CAM that localizes class-discriminative regions~\cite{selvaraju2017grad}, LIME that fits local linear models~\cite{ribeiro2016should}, and SHAP that provides game-theoretic feature attributions~\cite{lundberg2017unified}.

Perturbation-based methods offer a model-agnostic approach to attribution: by perturbing inputs or intermediate representations and observing changes in output, one can infer which features are important without access to model internals~\cite{zeiler2014visualizing}.
This approach has been applied to image classification by occluding image regions~\cite{zeiler2014visualizing}, to text analysis by removing words~\cite{li2016understanding}, and to tabular data by shuffling feature values~\cite{fisher2019all}.
Meaningful-perturbation methods~\cite{fong2017perturbation} and the removal-based view that unifies occlusion, permutation and Shapley-type attributions~\cite{covert2021removing} provide the general framework; a survey of perturbation-based explanation methods is given in~\cite{ivanovs2021perturbation}.

In the context of autoencoders, latent space perturbation provides insights into what the model has learned~\cite{higgins2017beta}.
By perturbing individual latent dimensions and observing reconstruction changes, one can identify which input features are encoded by each dimension.
This principle underlies disentanglement analysis in VAEs~\cite{chen2018isolating} and has been used for concept discovery in generative models~\cite{voynov2020unsupervised}.
The present work adapts latent perturbation for wavelength selection in ME-HSI: rather than seeking disentangled concepts, reconstruction sensitivity is traced back to specific excitation-emission combinations to identify informative wavelengths.

\subsection{Gap Addressed by This Work}
\label{sec:related_summary}

In summary, the band-selection families of Section~\ref{sec:related_band} (items 1--8) either rely on labels (items 2, 3, 6 and 8), fix the number of bands at training time (items 4 and 5), or evaluate bands along a single spectral axis (all items), and none accommodates a variable number of emission bands per excitation.
Deep learning approaches for hyperspectral imaging have focused primarily on supervised classification rather than unsupervised wavelength analysis.
Meanwhile, explainability methods have been developed for understanding model predictions but have not been systematically applied to spectral band selection.

The proposed framework addresses these gaps by: (1)~designing a 4D-aware autoencoder architecture with parallel excitation branches that accommodates variable band counts, (2)~leveraging unsupervised representation learning to avoid dependence on class labels, and (3)~applying perturbation-based attribution to trace learned importance back to interpretable wavelength selections.
This combination enables principled wavelength selection for ME-HSI data without requiring ground-truth annotations.